\subsection{Solution Output}
\label{SEC:Files:Solution_Output}

%%%%%%%%%%%%%%%%%%%%%%%%%%%%%%%%%%%%%%%%%%%%%%%%%%%%%%%%%%%%%%%%%%%%%%%%%%%%%%%%

At the start of every run job, CGNS grid and solution files are output.
%
%During a run job, a new CGNS solution file can be output at a specified
%iteration interval \andor time frequency.
% 
During a run job, if the current iteration is a multiple of
\inp[nmlOutput::OutputInterval]{Output\_Interval}
\andor the current simulation time is a multiple of
\inp[nmlOutput::OutputFrequency]{Output\_Frequency},
an updated CGNS solution file is output.
%
This update is based upon the current {\em aggregate} iteration or time
(including that from prior run jobs),
and not the {\em elapsed} value since the start of the current
run job or the last CGNS output.
%
The iteration number will be part of the CGNS solution file name.

%%%%%%%%%%%%%%%%%%%%%%%%%%%%%%%%%%%%%%%%%%%%%%%%%%%%%%%%%%%%%%%%%%%%%%%%%%%%%%%%

If the code terminates prematurely due to a detected NaN value,
a CGNS solution file is written for the previous time step.
%
The file name will include ``NaN'' to designate the solution was
unstable.
%
This file can be extremely helpful to help diagnose where the solution might
have diverged.

%%%%%%%%%%%%%%%%%%%%%%%%%%%%%%%%%%%%%%%%%%%%%%%%%%%%%%%%%%%%%%%%%%%%%%%%%%%%%%%%

At the end of every run job, an updated CGNS solution file is output
for that iteration.
%
A special CGNS file called \file{solution.time.cgns} is also output.
%
This file contains CGNS links to each of the instantaneous CGNS
solution files generated during the run, and can be used with flow
visualization software to easily read and examine all the time strand
solutions.
%
Beware that if writing to an existing directory, \file{solution.time.cgns}
will be overwritten with links to only those solution files generated
during the current run. 

%%%%%%%%%%%%%%%%%%%%%%%%%%%%%%%%%%%%%%%%%%%%%%%%%%%%%%%%%%%%%%%%%%%%%%%%%%%%%%%%

Every CGNS solution file contains an internal CGNS link to the
corresponding grid file.
%
For most plotting software, the user will not need to explicitly
load the corresponding grid file.
%
However, the grid file cannot be renamed without breaking the internal CGNS
link, and \emph{the grid file must be located within the same directory as the
solution file}.

%%%%%%%%%%%%%%%%%%%%%%%%%%%%%%%%%%%%%%%%%%%%%%%%%%%%%%%%%%%%%%%%%%%%%%%%%%%%%%%%

\NOTE
When running in parallel, the internal format of the output CGNS grid and
solution files may depend on the number of processors used.
%
If writing to an existing directory, then \file{solution.grid.cgns}
will overwriten.
%
This can cause problems if the number of processors has changed between
successive run jobs, because each previous solution file has an internal
link to \file{solution.grid.cgns} and the internal format of that file
may no longer be compatible with the old solution file.
%
Therefore, it is recommended that each successive run job use a different
output directory.
%
This way, the {\tt solution.grid.cgns} and {\tt solution.time.cgns} files
for each run will be preserved.

%%%%%%%%%%%%%%%%%%%%%%%%%%%%%%%%%%%%%%%%%%%%%%%%%%%%%%%%%%%%%%%%%%%%%%%%%%%%%%%%

\IDEA{Add information about CGNS output being slow for large numbers of MPI
processes, so it is suggested to disable CGNS output altogether and to use the
GFR restart-to-CGNS utility to convert restart files to CGNS solution files.}

%%%%%%%%%%%%%%%%%%%%%%%%%%%%%%%%%%%%%%%%%%%%%%%%%%%%%%%%%%%%%%%%%%%%%%%%%%%%%%%%

\SeeAlso{
\inp[nmlOutput::OutputDir]{Output\_Dir},
\inp[nmlOutput::OutputInterval]{Output\_Interval},
\inp[nmlOutput::OutputFrequency]{Output\_Frequency}
}

%%%%%%%%%%%%%%%%%%%%%%%%%%%%%%%%%%%%%%%%%%%%%%%%%%%%%%%%%%%%%%%%%%%%%%%%%%%%%%%%

